\chapter{Fonctionnalités principales}

Ce chapitre fixe le périmètre de la version rendue en fin de trimestre. Chaque
exigence porte un identifiant et une priorité :

\begin{description}
  \item[Haute] indispensable à la démonstration ;
  \item[Moyenne] prévue, mais peut être simplifiée si le temps manque ;
  \item[Basse] finition, traitée en dernier.
\end{description}

Les numéros d'écran renvoient aux maquettes du chapitre~\ref{ch:interface}.

\section{Premier lancement et préférences}

Au premier lancement, quatre écrans courts recueillent ce qui sert ensuite à
personnaliser l'application (écrans 6 à 10, figure~\ref{fig:lancement}).
Les choix restent modifiables dans les réglages.

\begin{exigence}[priorite=haute]{Vérification d'âge}
  L'utilisateur déclare s'il a 18 ans ou plus. S'il répond non, ou s'il choisit
  de continuer sans alcool, les vins et cocktails alcoolisés sont masqués dans
  toute l'application. La déclaration suffit, aucun justificatif n'est demandé.
  Chaque écran qui affiche de l'alcool porte la mention sanitaire légale.
\end{exigence}

\begin{exigence}[priorite=haute]{Régime alimentaire}
  Choix multiple parmi : de tout, végétarien, végétalien, pescétarien, sans
  gluten, sans lactose. L'étape peut être passée.
\end{exigence}

\begin{exigence}[priorite=moyenne]{Matériel de cuisine}
  L'utilisateur coche ses appareils : plaques, four, micro-ondes, air fryer,
  mixeur, grill. Les recettes qui en demandent d'autres sont signalées.
\end{exigence}

\begin{exigence}[priorite=haute]{Boissons}
  Deux options : vins et cocktails, ou uniquement sans alcool. Ce choix
  détermine les accords proposés partout dans l'application.
\end{exigence}

\begin{exigence}[priorite=haute]{Recettes marquées selon le profil}
  Les préférences ne masquent pas les recettes : elles les marquent. Une
  recette est compatible, adaptable (avec la substitution proposée, par
  exemple des pâtes sans gluten) ou signalée. Les filtres de l'onglet
  Explorer sont préactivés selon le profil. Les informations sur les
  allergènes restent indicatives.
\end{exigence}

\section{Carte du jour}

L'accueil présente le repas du jour sous la forme d'une carte de restaurant
(écrans 1 et 2).

\begin{exigence}[priorite=haute]{Repas du jour}
  Chaque jour, l'accueil propose un plat, un verre et un dessert compatibles
  avec le profil. Le tirage est enregistré avec la date : la carte reste la
  même toute la journée, y compris hors ligne.
\end{exigence}

\begin{exigence}[priorite=moyenne]{Raccourcis de découverte}
  Sous la carte, des raccourcis par envie (pâtes, mer, végétarien,
  desserts…) et par pays ouvrent l'onglet Explorer avec le filtre
  correspondant.
\end{exigence}

\section{Recherche et filtrage}

L'onglet Explorer regroupe les plats, les cocktails et les vins
(écrans 11, 12, 14 et 20, figures~\ref{fig:onglets} et~\ref{fig:exploration}).

\begin{exigence}[priorite=haute]{Recherche}
  Un champ unique cherche par nom de plat, par ingrédient principal ou par
  pays. Les résultats s'affichent dans l'onglet actif : plats, cocktails ou
  vins.
\end{exigence}

\begin{exigence}[priorite=haute]{Filtres des plats}
  Catégorie, pays, régime et matériel. Les filtres actifs restent visibles
  au-dessus des résultats et se retirent d'un toucher.
\end{exigence}

\begin{exigence}[priorite=moyenne]{Filtres des cocktails}
  Avec ou sans alcool, et type de verre (highball, coupe, tumbler…).
\end{exigence}

\begin{exigence}[priorite=moyenne]{Catalogue des vins}
  Vins classés par couleur (rouge, blanc, rosé, effervescent). Chaque cépage
  indique son profil, sa température de service et les plats avec lesquels
  il s'accorde.
\end{exigence}

\begin{exigence}[priorite=moyenne]{Carte du monde}
  Liste des pays regroupés par région, chacun avec un plat emblématique.
  Toucher un pays ouvre toutes ses recettes.
\end{exigence}

\section{Fiches recettes et mode cuisine}

Écrans 3, 4, 13 et 15, figures~\ref{fig:recette} et~\ref{fig:exploration}.

\begin{exigence}[priorite=haute]{Fiche recette}
  Photo, pays, catégorie, ingrédients présentés en étagère, étapes de
  préparation et badge d'adaptation au profil. Le matériel nécessaire et
  les durées s'affichent quand ils sont repérés dans les instructions. Deux
  actions : cuisiner et ajouter à la liste de courses.
\end{exigence}

\begin{exigence}[priorite=moyenne]{Ajustement des portions}
  Un compteur modifie le nombre de personnes. Les quantités numériques sont
  recalculées ; les autres (« au goût », « une pincée ») restent telles
  quelles. TheMealDB ne précisant pas le nombre de portions, une recette est
  considérée comme prévue pour quatre personnes.
\end{exigence}

\begin{exigence}[priorite=moyenne]{Ingrédients à cocher}
  L'utilisateur coche ce qu'il a déjà. Seuls les ingrédients non cochés
  partent dans la liste de courses.
\end{exigence}

\begin{exigence}[priorite=haute]{Mode cuisine}
  Une étape par écran, en grand, avec les ingrédients qu'elle utilise et une
  barre de progression. Quand une étape mentionne une durée, un minuteur est
  proposé ; il continue de tourner d'une étape à l'autre. L'écran reste
  allumé pendant toute la préparation.
\end{exigence}

\begin{exigence}[priorite=haute]{Fiche cocktail}
  Photo, ingrédients et doses, verre recommandé, mention avec ou sans alcool,
  et instructions.
\end{exigence}

\section{Accords plats-boissons et menus}

Écrans 5 et 16, figures~\ref{fig:onglets} et~\ref{fig:menu-scan}.

\begin{exigence}[priorite=haute]{La table dressée}
  Pour un plat donné, l'onglet Accords propose un vin, un cocktail et un
  dessert, disposés comme un couvert vu du dessus. Chaque élément ouvre sa
  fiche. Un sélecteur bascule entre accords avec et sans alcool.
\end{exigence}

\begin{exigence}[priorite=haute]{Règles d'accord}
  Les accords suivent une table de règles rédigée par l'équipe, qui associe
  la catégorie et le pays d'un plat à des styles de vin et de cocktail (par
  exemple, fruits de mer et vin blanc sec). Aucune des deux API ne fournit
  cette information.
\end{exigence}

\begin{exigence}[priorite=haute]{Composer un menu}
  Un menu réunit une entrée, un plat, un dessert et une boisson, pour une
  occasion et un nombre de convives. Il est enregistré, et ses ingrédients
  s'ajoutent à la liste de courses en une fois, ajustés au nombre de
  convives.
\end{exigence}

\begin{exigence}[priorite=basse]{Partager un menu}
  Le menu est envoyé en texte par la feuille de partage d'Android.
\end{exigence}

\section{Liste de courses}

Écran 17, figure~\ref{fig:onglets}. La liste occupe un onglet de la
navigation et ne dépend d'aucun service externe : elle fait donc partie de la
version de base.

\begin{exigence}[priorite=haute]{Ajout depuis une recette ou un menu}
  Les ingrédients arrivent avec les quantités ajustées. Un même ingrédient
  exprimé dans la même unité est additionné ; sinon, les deux lignes
  restent séparées.
\end{exigence}

\begin{exigence}[priorite=moyenne]{Classement par rayon}
  Les articles sont regroupés par rayon (fruits et légumes, crèmerie,
  épicerie, cave…), avec l'avancement de chaque rayon. Un ingrédient inconnu
  va dans « Autres ».
\end{exigence}

\begin{exigence}[priorite=haute]{Cocher en magasin}
  Un toucher coche l'article. Le texte n'est jamais barré, pour rester lisible.
\end{exigence}

\section{Favoris}

\begin{exigence}[priorite=haute]{Favoris}
  Plats, boissons et menus se mettent en favori d'un toucher sur le cœur. Ils
  sont rangés dans trois onglets et restent disponibles hors ligne
  (écran 19).
\end{exigence}

\section{Consultation hors ligne}

\begin{exigence}[priorite=haute]{Données disponibles sans réseau}
  Sans connexion, l'utilisateur garde l'accès à la carte du jour, à ses
  favoris, à ses menus, à sa liste de courses et aux dernières fiches
  consultées, photos comprises. La recherche nécessite le réseau ; hors
  ligne, un message l'indique clairement au lieu d'une liste vide.
\end{exigence}

\section{Réglages et profil}

\begin{exigence}[priorite=moyenne]{Réglages}
  Modifier le régime, le matériel et le mode sans alcool ; choisir le thème
  (système, clair, sombre) ; activer le retour haptique. Les animations sont
  réduites automatiquement si le téléphone le demande. Une page À propos
  cite les sources, les mentions légales et les licences open source
  (écrans 21 et 22).
\end{exigence}

\begin{exigence}[priorite=basse]{Profil local}
  L'icône de profil ouvre une connexion simulée : un prénom et un avatar
  enregistrés sur le téléphone. Il n'y a ni compte ni authentification.
\end{exigence}
